\documentclass[11pt]{article}
\newcommand{\ListaTitulo}{Lista 10 --- Friedman e Blocos Incompletos}
% Preâmbulo compartilhado — MATD48, Listas de Exercícios 2026
% Usado por Lista{NN}.tex e Gabarito{NN}.tex via % Preâmbulo compartilhado — MATD48, Listas de Exercícios 2026
% Usado por Lista{NN}.tex e Gabarito{NN}.tex via % Preâmbulo compartilhado — MATD48, Listas de Exercícios 2026
% Usado por Lista{NN}.tex e Gabarito{NN}.tex via \input{preamble}

\usepackage[utf8]{inputenc}
\usepackage[T1]{fontenc}
\usepackage[brazil]{babel}
\usepackage{fullpage}
\usepackage{amsmath,amssymb}
\usepackage{enumitem}
\usepackage{xcolor}
\usepackage{fancyhdr}
\usepackage{titlesec}
\usepackage{listings}
\usepackage{hyperref}
\usepackage{booktabs}
\usepackage{graphicx}

\definecolor{matd48azul}{HTML}{0B4F6C}
\definecolor{matd48verde}{HTML}{2E8B57}

\titleformat{\section}{\normalfont\Large\bfseries\color{matd48azul}}{\thesection}{1em}{}
\titleformat{\subsection}{\normalfont\large\bfseries\color{matd48azul}}{\thesubsection}{1em}{}

\providecommand{\ListaTitulo}{Lista de Exercícios}

\setlength{\headheight}{14pt}
\pagestyle{fancy}
\fancyhf{}
\lhead{\small MATD48 --- Planejamento de Experimentos A}
\rhead{\small \ListaTitulo}
\cfoot{\thepage}
\renewcommand{\headrulewidth}{0.4pt}

\lstset{
  basicstyle=\ttfamily\small,
  breaklines=true,
  frame=single,
  columns=fullflexible,
  backgroundcolor=\color{gray!8},
  commentstyle=\color{matd48verde},
  keywordstyle=\color{matd48azul}\bfseries,
  language=R,
  extendedchars=true,
  literate=%
    {á}{{\'a}}1 {à}{{\`a}}1 {â}{{\^a}}1 {ã}{{\~a}}1
    {é}{{\'e}}1 {ê}{{\^e}}1 {í}{{\'i}}1 {ó}{{\'o}}1
    {ô}{{\^o}}1 {õ}{{\~o}}1 {ú}{{\'u}}1 {ç}{{\c c}}1
    {Á}{{\'A}}1 {À}{{\`A}}1 {Â}{{\^A}}1 {Ã}{{\~A}}1
    {É}{{\'E}}1 {Ê}{{\^E}}1 {Í}{{\'I}}1 {Ó}{{\'O}}1
    {Ô}{{\^O}}1 {Õ}{{\~O}}1 {Ú}{{\'U}}1 {Ç}{{\c C}}1
}

\hypersetup{colorlinks=true, linkcolor=matd48azul, urlcolor=matd48azul, citecolor=matd48azul}

\newcommand{\cabecalho}[2]{%
  \begin{center}
    {\Large\bfseries\color{matd48azul} #1}\\[0.3em]
    {\large #2}\\[0.3em]
    \rule{\linewidth}{0.6pt}
  \end{center}
  \vspace{0.5em}
}

\newenvironment{questao}[1]{\vspace{1em}\noindent\textbf{Questão #1.}\hspace{0.4em}}{\par}
\newenvironment{solucao}{\vspace{0.3em}\noindent\textit{Solução.}\hspace{0.4em}\color{black!85}}{\par\vspace{0.5em}}


\usepackage[utf8]{inputenc}
\usepackage[T1]{fontenc}
\usepackage[brazil]{babel}
\usepackage{fullpage}
\usepackage{amsmath,amssymb}
\usepackage{enumitem}
\usepackage{xcolor}
\usepackage{fancyhdr}
\usepackage{titlesec}
\usepackage{listings}
\usepackage{hyperref}
\usepackage{booktabs}
\usepackage{graphicx}

\definecolor{matd48azul}{HTML}{0B4F6C}
\definecolor{matd48verde}{HTML}{2E8B57}

\titleformat{\section}{\normalfont\Large\bfseries\color{matd48azul}}{\thesection}{1em}{}
\titleformat{\subsection}{\normalfont\large\bfseries\color{matd48azul}}{\thesubsection}{1em}{}

\providecommand{\ListaTitulo}{Lista de Exercícios}

\setlength{\headheight}{14pt}
\pagestyle{fancy}
\fancyhf{}
\lhead{\small MATD48 --- Planejamento de Experimentos A}
\rhead{\small \ListaTitulo}
\cfoot{\thepage}
\renewcommand{\headrulewidth}{0.4pt}

\lstset{
  basicstyle=\ttfamily\small,
  breaklines=true,
  frame=single,
  columns=fullflexible,
  backgroundcolor=\color{gray!8},
  commentstyle=\color{matd48verde},
  keywordstyle=\color{matd48azul}\bfseries,
  language=R,
  extendedchars=true,
  literate=%
    {á}{{\'a}}1 {à}{{\`a}}1 {â}{{\^a}}1 {ã}{{\~a}}1
    {é}{{\'e}}1 {ê}{{\^e}}1 {í}{{\'i}}1 {ó}{{\'o}}1
    {ô}{{\^o}}1 {õ}{{\~o}}1 {ú}{{\'u}}1 {ç}{{\c c}}1
    {Á}{{\'A}}1 {À}{{\`A}}1 {Â}{{\^A}}1 {Ã}{{\~A}}1
    {É}{{\'E}}1 {Ê}{{\^E}}1 {Í}{{\'I}}1 {Ó}{{\'O}}1
    {Ô}{{\^O}}1 {Õ}{{\~O}}1 {Ú}{{\'U}}1 {Ç}{{\c C}}1
}

\hypersetup{colorlinks=true, linkcolor=matd48azul, urlcolor=matd48azul, citecolor=matd48azul}

\newcommand{\cabecalho}[2]{%
  \begin{center}
    {\Large\bfseries\color{matd48azul} #1}\\[0.3em]
    {\large #2}\\[0.3em]
    \rule{\linewidth}{0.6pt}
  \end{center}
  \vspace{0.5em}
}

\newenvironment{questao}[1]{\vspace{1em}\noindent\textbf{Questão #1.}\hspace{0.4em}}{\par}
\newenvironment{solucao}{\vspace{0.3em}\noindent\textit{Solução.}\hspace{0.4em}\color{black!85}}{\par\vspace{0.5em}}


\usepackage[utf8]{inputenc}
\usepackage[T1]{fontenc}
\usepackage[brazil]{babel}
\usepackage{fullpage}
\usepackage{amsmath,amssymb}
\usepackage{enumitem}
\usepackage{xcolor}
\usepackage{fancyhdr}
\usepackage{titlesec}
\usepackage{listings}
\usepackage{hyperref}
\usepackage{booktabs}
\usepackage{graphicx}

\definecolor{matd48azul}{HTML}{0B4F6C}
\definecolor{matd48verde}{HTML}{2E8B57}

\titleformat{\section}{\normalfont\Large\bfseries\color{matd48azul}}{\thesection}{1em}{}
\titleformat{\subsection}{\normalfont\large\bfseries\color{matd48azul}}{\thesubsection}{1em}{}

\providecommand{\ListaTitulo}{Lista de Exercícios}

\setlength{\headheight}{14pt}
\pagestyle{fancy}
\fancyhf{}
\lhead{\small MATD48 --- Planejamento de Experimentos A}
\rhead{\small \ListaTitulo}
\cfoot{\thepage}
\renewcommand{\headrulewidth}{0.4pt}

\lstset{
  basicstyle=\ttfamily\small,
  breaklines=true,
  frame=single,
  columns=fullflexible,
  backgroundcolor=\color{gray!8},
  commentstyle=\color{matd48verde},
  keywordstyle=\color{matd48azul}\bfseries,
  language=R,
  extendedchars=true,
  literate=%
    {á}{{\'a}}1 {à}{{\`a}}1 {â}{{\^a}}1 {ã}{{\~a}}1
    {é}{{\'e}}1 {ê}{{\^e}}1 {í}{{\'i}}1 {ó}{{\'o}}1
    {ô}{{\^o}}1 {õ}{{\~o}}1 {ú}{{\'u}}1 {ç}{{\c c}}1
    {Á}{{\'A}}1 {À}{{\`A}}1 {Â}{{\^A}}1 {Ã}{{\~A}}1
    {É}{{\'E}}1 {Ê}{{\^E}}1 {Í}{{\'I}}1 {Ó}{{\'O}}1
    {Ô}{{\^O}}1 {Õ}{{\~O}}1 {Ú}{{\'U}}1 {Ç}{{\c C}}1
}

\hypersetup{colorlinks=true, linkcolor=matd48azul, urlcolor=matd48azul, citecolor=matd48azul}

\newcommand{\cabecalho}[2]{%
  \begin{center}
    {\Large\bfseries\color{matd48azul} #1}\\[0.3em]
    {\large #2}\\[0.3em]
    \rule{\linewidth}{0.6pt}
  \end{center}
  \vspace{0.5em}
}

\newenvironment{questao}[1]{\vspace{1em}\noindent\textbf{Questão #1.}\hspace{0.4em}}{\par}
\newenvironment{solucao}{\vspace{0.3em}\noindent\textit{Solução.}\hspace{0.4em}\color{black!85}}{\par\vspace{0.5em}}


\begin{document}

\cabecalho{Lista de Exercícios 10}{Teste de Friedman e blocos incompletos balanceados (Capítulo 4 / Aula 10)}

\noindent \textit{Entrega: até a próxima aula. Justifique todas as respostas — nestas listas, a
justificativa vale mais que a resposta final. O gabarito completo está em \texttt{Gabarito10.pdf}.}

\begin{questao}{1} \textbf{(Psicologia --- estatística de Friedman)}

Um estudo de medidas repetidas compara $t=3$ técnicas de estudo em $b=10$ estudantes (o bloco).
Dentro de cada estudante, as três técnicas são ranqueadas de 1 (pior) a 3 (melhor) quanto à nota
obtida. As somas de posto por técnica foram $R_1=12$, $R_2=20$, $R_3=28$.

\begin{enumerate}[label=(\alph*)]
  \item Verifique que $R_1+R_2+R_3$ é consistente com $b=10$ estudantes ranqueando 3 técnicas
    cada (qual deve ser essa soma total, e por quê?).
  \item Calcule a estatística de Friedman $F_r$.
  \item Sabendo que $\chi^2_{0{,}05;\,2}=5{,}99$, há evidência de diferença sistemática entre as
    técnicas? Justifique.
\end{enumerate}
\end{questao}

\begin{questao}{2} \textbf{(Dedução --- valor esperado da soma de postos sob $H_0$)}

Sob $H_0$ do teste de Friedman, dentro de cada bloco todas as $t!$ ordenações dos $t$
tratamentos são igualmente prováveis.

\begin{enumerate}[label=(\alph*)]
  \item Argumente, por simetria, que sob $H_0$ o posto atribuído a um tratamento específico
    \emph{dentro de um bloco} é uniformemente distribuído sobre $\{1,2,\dots,t\}$.
  \item Use o item (a) para mostrar que $\mathbb E[\text{posto}(Y_{ij})] = \frac{t+1}{2}$ para
    todo tratamento $i$ e todo bloco $j$, sob $H_0$.
  \item Conclua que $\mathbb E[R_i] = \frac{b(t+1)}{2}$ para todo tratamento $i$, sob $H_0$, e
    confira que essa é a média (não necessariamente o valor) que cada $R_i$ do exercício anterior
    deveria ter, em expectativa, se as três técnicas fossem de fato equivalentes.
\end{enumerate}
\end{questao}

\begin{questao}{3} \textbf{(Agricultura --- parâmetros de um BIB)}

Uma cooperativa quer montar um painel de degustação com $t=6$ variedades de café, mas cada
degustador (bloco) só avalia $k=3$ xícaras por sessão.

\begin{enumerate}[label=(\alph*)]
  \item Usando as identidades $bk=tr$ e $\lambda(t-1)=r(k-1)$, encontre o menor valor inteiro
    positivo de $\lambda$ para o qual $r$ resulta inteiro, e calcule o $r$ e o $b$
    correspondentes.
  \item Verifique numericamente as duas identidades com os valores encontrados.
\end{enumerate}
\end{questao}

\begin{questao}{4} \textbf{(Agricultura --- verificando e completando um BIB)}

Para $t=4$ variedades de tomate e blocos de tamanho $k=3$, um técnico propõe inicialmente estes
três blocos:
$$
\text{Bloco 1}=\{1,2,3\}, \quad \text{Bloco 2}=\{1,2,4\}, \quad \text{Bloco 3}=\{1,3,4\}.
$$

\begin{enumerate}[label=(\alph*)]
  \item Calcule $r_i$ (número de blocos em que aparece) para cada uma das 4 variedades. Esse
    arranjo de 3 blocos já é um BIB válido? Justifique por que a condição de $r$ constante entre
    tratamentos é necessária (ainda que não suficiente) para balanceamento.
  \item Proponha um quarto bloco que corrija o problema identificado em (a), e verifique, para o
    arranjo de 4 blocos resultante, que todo par de variedades coocorre o mesmo número $\lambda$
    de vezes.
\end{enumerate}
\end{questao}

\begin{questao}{5} \textbf{(Dissertativa --- Friedman vs. DBCA, blocos completos vs. incompletos)}

\begin{enumerate}[label=(\alph*)]
  \item Um pesquisador tem uma resposta claramente normal, sem outliers, com $b=20$ blocos. Ele
    ainda assim decide usar o teste de Friedman "por segurança". Que preço estatístico ele paga
    por essa escolha, mesmo que a conclusão qualitativa não mude?
  \item Em que situação um pesquisador deveria preferir um delineamento em blocos
    \emph{incompletos} balanceados a um DBCA, mesmo quando seria fisicamente possível caber todos
    os tratamentos em cada bloco?
\end{enumerate}
\end{questao}

\end{document}
